\begin{abstract}
The semiring framework and its extensions form the basis of a rich collection of theoretical
results and implementations for provenance tracking of database queries.
Many real-world queries use aggregation and conditions on the aggregate
values.
Support for such queries has been proposed by introducing semimodule
elements as aggregate values and formal comparisons between aggregate values as tuple
annotations, which takes the approach outside the standard semiring framework. In this work, we
show how to introduce a semantics for the provenance of such queries in
arbitrary commutative semirings with monus (or \emph{m-semirings}), without the need for additional
operators. This semantics is shown to agree with the standard provenance of the
aggregation-free self-join rewriting of \sql{HAVING COUNT(*)} queries in
semirings that are absorptive and where times distributes over monus. We
derive algorithms for this semantics and implement them within the
ProvSQL system, with viable performance on a real-world dataset for
probabilistic query evaluation.
\end{abstract}
